% =====================================================================
% Section: Experiments -- Datasets and implementation details (outline.md §3.1;
%   defines sec:exp_setup). Opens the Experiments section.
% Draft v1 (2026-06-17). All settings are transcribed verbatim from the training
%   configs (CATANet/experiments/*/*.yml) and the efficiency/test logs recorded in
%   paper/plan/data-collection-checklist.md (A2/A3) and experiment-protocol.md.
% This draft makes NO result assertions; it only states the protocol.
% HARDWARE PLACEHOLDER: the exact GPU model/count and total GPU-hours are not yet
%   collected (data-collection-checklist A1, still open). The bracketed
%   [PLACEHOLDER ...] must be filled from the training machine before submission;
%   do NOT fabricate a specific device.
% =====================================================================
\section{Experiments}
\label{sec:experiments}

\subsection{Datasets and implementation details}
\label{sec:exp_setup}

\paragraph{Datasets and metrics.}
Following the standard lightweight-SR protocol, we train on the $800$ training
images of DIV2K~\cite{agustsson2017div2k} and evaluate on five benchmarks: Set5,
Set14, BSD100, Urban100, and Manga109. Low-resolution inputs are produced by
bicubic downscaling at $\times2$, $\times3$, and $\times4$. We report PSNR and SSIM
on the Y channel of the YCbCr space, discarding a border of $s$ pixels for scale
$s$, which matches the convention used by the cited baselines.

\paragraph{Network configuration.}
DPRNet keeps the CATANet backbone unchanged and replaces only the routing inside
each TAB with the DPR module (Sec.~\ref{sec:method}). The deep feature extractor
stacks $L{=}8$ blocks with channel width $C{=}40$. The number of prototype slots
per block is $M=[16,32,64,128,16,32,64,128]$, and the intra-group attention uses a
group size of $128$. The learnable routing temperature is initialized at
$e^{6}$ via a log-scale and clamped at $\tau_{\max}{=}10$
(Eq.~\eqref{eq:scores}); the per-block usage-balancing weights are
$\lambda=[5{\times}10^{-4}]_{\times6}$ for the first six blocks and
$7{\times}10^{-4}$ for the last two (Eq.~\eqref{eq:balance}). These yield
$0.60$--$0.67$\,M parameters depending on scale (Table~\ref{tab:efficiency}).

\paragraph{Training.}
We first train the $\times2$ model from scratch for $800$k iterations, and then
finetune it to $\times3$ and $\times4$ for $250$k iterations each, initializing from
the $\times2$ checkpoint. Training patches are $128{\times}128$, $192{\times}192$,
and $256{\times}256$ HR crops for $\times2/\times3/\times4$ respectively, with a
batch size of $16$ and random horizontal flips and $90^{\circ}$ rotations for
augmentation. We optimize the L1 reconstruction loss together with the
usage-balancing regularizer (Eq.~\eqref{eq:balance}) using Adam
($\beta_1{=}0.9$, $\beta_2{=}0.99$, no weight decay) at an initial learning rate of
$2{\times}10^{-4}$, halved by a MultiStep schedule (milestones at
$\{300,500,650,700,750\}$k for $\times2$ and $\{125,200,225,237.5\}$k for the
$\times3/\times4$ finetuning). The learnable routing temperature is trained with a
$0.1{\times}$ learning-rate multiplier for stability. Validation is run every
$2{.}5$k iterations; the unified reporting checkpoint per scale is selected by the
best benchmark-averaged PSNR/SSIM (Set5--Manga109) without per-dataset
cherry-picking. Experiments use PyTorch with the BasicSR
framework~\cite{wang2020basicsr} on [PLACEHOLDER: $N\times$ NVIDIA \texttt{<GPU
model>}, total $\approx$\texttt{<GPU-hours>}; fill from training machine before
submission].

\paragraph{Efficiency measurement.}
Parameters, Multi-Adds (MACs reported by thop), and inference latency are measured
on a single CUDA GPU for a fixed HR output of $\approx720{\times}1280$ (the LR input
is derived per scale), averaging over $100$ runs after $20$ warm-up runs. To compare
with FLOPs-based reports, the Multi-Adds values are multiplied by two.
